(1) 新阱 $0<x<2L$ 的归一化本征态为 $\phi_n(x)=L^{-1/2}\sin(n\pi x/(2L))$，$n\ge1$，能级 $E'_n=n^2\pi^2\hbar^2/(8mL^2)$。量子态在瞬时改变势壁时连续，但原波函数在 $L<x<2L$ 处初始为零，不会瞬间变成新阱基态。
(2) $c_n=\int_0^{2L}\phi_n^*(x)\psi_i(x)dx=(\sqrt2/L)\int_0^L\sin(\pi x/L)\sin(n\pi x/(2L))dx$。当 $n\ne2$，积分得 $c_n=(\sqrt2/\pi)\sin(n\pi/2)/(1-n^2/4)$；$n=2$ 是可去奇点，直接积分得 $c_2=1/\sqrt2$。其余偶数 $n\ne2$ 的系数为零。
(3) 新阱第二能级测量概率 $P_2=|c_2|^2=1/2$。瞬态波函数在 $x=L$ 连续且为零，故动能期望的二次型积分不变：$\langle H_{\rm new}\rangle=(\hbar^2/(2m))\int_0^L|\psi_i'(x)|^2dx=\pi^2\hbar^2/(2mL^2)$，亦即原阱基态能量；不可以用最可能的某个新能级替代平均值。
(4) 绝热慢移右壁使非简并基态追随到宽度 $2L$ 的新基态，其末态能量是 $E'_1=\pi^2\hbar^2/(8mL^2)$，为原能量四分之一；理想瞬移在瞬间未改变动能期望。完备性保证 $\sum_{n\ge1}|c_n|^2=1$，其成立依赖包含 $n=2$ 的特别积分值与初态在新区域的零延拓。
